\section{Introduction}
\label{sec:intro}

Dense visual SLAM has converged on 3D Gaussian splatting~\cite{kerbl20233dgs}
as its map representation, because Gaussians render photorealistically in real
time where neural fields~\cite{mildenhall2020nerf} do not. The dominant
designs --- MonoGS~\cite{matsuki2024monogs},
Photo-SLAM~\cite{huang2024photoslam}, SplaTAM~\cite{keetha2024splatam} ---
obtain their Gaussians by \emph{per-scene optimisation}: primitives are
initialised, then densified and refined by gradient descent against the
incoming frames.

A different route has become available. Feed-forward reconstruction networks
predict geometry directly from images without optimisation:
DUSt3R~\cite{wang2024dust3r} regresses pointmaps from an image pair,
MASt3R~\cite{leroy2024mast3r} adds matching, and
Splatt3R~\cite{smart2024splatt3r} extends MASt3R with a head that emits a full
Gaussian per pixel --- position, covariance, colour and opacity --- in a single
forward pass. MASt3R-SLAM~\cite{murai2025mast3rslam} shows that the MASt3R
prior supports real-time dense tracking. The combination is therefore
attractive: a SLAM system whose map appears at the speed of a forward pass.

The difficulty is a mismatch of regimes. Splatt3R is trained on
\emph{curated image pairs}, and its loss sees exactly two views at a time. A
SLAM system does something categorically different: it fuses hundreds of
keyframes onto the same physical surfaces, so every surface is covered many
times over, by predictions made from different viewpoints, at different
exposures, with different confidences. Nothing in pairwise training observes
this crowding. Our starting observation is that a checkpoint which is
excellent per pair can nonetheless produce a fused map that is hazy and
over-solid --- and that fixing this is not a matter of training longer.

This paper asks what actually has to change, and reports the answer as three
results.

\paragraph{Where adaptation is admissible.}
The obvious move --- adapt the backbone to the target domain --- fails
decisively. Encoder LoRA~\cite{hu2022lora} degrades reconstruction by $49\%$
with a $1821\times$ scale explosion, and decoder-only LoRA improves for one
epoch then collapses. The reason is structural: Splatt3R's Gaussian head was
trained against \emph{frozen} encoder features, and perturbing those features
feeds the head inputs it has never seen. Training only the Gaussian head, with
the encoder bit-frozen, improves deployed map quality across five dataset
families. Because the encoder is untouched, tracking is unchanged by
construction --- a property we verify externally rather than assume.

\paragraph{A training-free lever, and why training cannot replace it.}
Fading each Gaussian's opacity at the moment it is injected into the map, in
proportion to how little the pointmap vouches for it, improves the online map
in $10$ of $12$ cells across four families, six heads and eight scenes, with
two ties and \emph{no} harmful cell. This is not a heuristic that a better
loss would subsume: crowding is a map-level property that only exists after
many keyframes accumulate, and a pair-training objective never sees a map. We
support this with a dose--response experiment at two thinning depths.

\paragraph{What opacity is.}
Both preceding results turn out to be the same result. We show that opacity in
this architecture is a \emph{brightness-trim dial} against the
black-background rasteriser, not the confidence or hedging channel it is
usually assumed to be: because the colour term is anchored to the (possibly
degraded) input pixel and is sigmoid-bounded, opacity is the only channel with
headroom to brighten. The decisive evidence is a pre-registered opposite-sign
prediction --- heads trained under synthetic darkening should correlate opacity
\emph{negatively} with input luminance, real-sensor heads \emph{positively} ---
which holds on all five families. Under this account the base checkpoint's
collapsed opacity means it \emph{cannot} trim; injection-time thinning supplies
the trim externally; and heads that deploy well are the ones that learned it.

\paragraph{An honest comparison, and what it cost us.}
We built MASt3R-SLAM, MonoGS, Photo-SLAM and VGGT-SLAM~\cite{maggio2025vggtslam}
from source and evaluated every system --- ours included --- through a single
protocol, on held-out views that are keyframes of no system. Two findings
emerged that we did not set out to make.

First, \emph{published GS-SLAM rendering numbers are not comparable to ours},
and the gap is quantifiable. Running Photo-SLAM ourselves on Replica office0
gives $22.2$\dB{} where third-party tables report ${\sim}30.9$\dB{} for the
same system. We independently reproduce an $8.8$\dB{} gap on MonoGS between
rendering its own map from its own estimated poses versus ground-truth poses.
Two systems, two methods, the same ${\sim}8.7$\dB{}: the literature and this
paper differ by one protocol's worth of evaluation choices.

Second, several of our own claims weakened when pushed. A three-scene Replica
subset suggested a $+5.29$\dB{} advantage; completing all eight scenes reduced
it to $+3.44$\dB{} and revealed that we win LPIPS on only $5/8$. On TUM at
matched compute we \emph{tie} MonoGS on PSNR and \emph{lose} on LPIPS. And our
maps carry $25$--$100\times$ more Gaussians at $8$--$16\times$ the peak memory;
pruned to a baseline's budget our quality falls $8.7$--$10.3$\dB{}
\emph{below} it, so the size is constitutive rather than redundant. We report
these as measured results, with the curve, rather than as caveats.

Our contributions are: (i) a head-only adaptation recipe for feed-forward
Gaussian SLAM, with the negative encoder-adaptation result that motivates it;
(ii) an injection-time thinning lever with a dose--response argument for why
training cannot substitute for it; (iii) a falsifiable mechanism for opacity,
confirmed by a pre-registered opposite-sign prediction on $5/5$ families;
(iv) a single-protocol comparison against four locally-built systems, with a
quantified protocol offset that we believe affects how this literature should
be read; and (v) a measured account of the efficiency deficit our
representation carries.
